%% =========================================================================
%% When Does LLM Unlearning Fail? Predicting and Protecting Susceptible Retained Behavior
%% From Update-Conditioned Susceptibility to Profile-Guided Protection
%%
%% KDD 2027 Research Track submission draft (San Jose, CA, August 2027).
%% Template: ACM acmart, sigconf + anonymous + review
%% (CFP-recommended setting; review adds line numbers for reviewers).
%% Build: pdflatex -> bibtex -> pdflatex x2  (Overleaf handles this).
%%
%% FILE STRUCTURE
%%   main.tex           preamble, metadata, and \input skeleton (this file)
%%   sections/*.tex     one file per section, numbered in reading order
%%   references.bib     bibliography
%%
%% The paper asks whether a pre-unlearning susceptibility profile predicts
%% realized collateral damage and improves protection of retained behavior.
%% =========================================================================
\documentclass[sigconf,anonymous,review]{acmart}

\AtBeginDocument{%
  \providecommand\BibTeX{{Bib\TeX}}}

%% ---- rights block: placeholder values; replaced at camera-ready.
\setcopyright{none}
\acmConference[KDD '27]{the 33rd ACM SIGKDD Conference on
  Knowledge Discovery and Data Mining}{August 2027}{San Jose, CA, USA}
\acmBooktitle{Proceedings of the 33rd ACM SIGKDD Conference on Knowledge
  Discovery and Data Mining (KDD '27), August 2027, San Jose, CA, USA}
\acmYear{2027}
\acmDOI{}
\acmISBN{}
%% \acmSubmissionID{TODO-openreview-id}

\usepackage{booktabs}
\usepackage{amsmath}
\usepackage{xcolor}
\usepackage{placeins}
\usepackage[ruled,linesnumbered]{algorithm2e}
\usepackage{tikz}
\usetikzlibrary{arrows.meta,calc,positioning,matrix,shapes}

%% ---- notation ------------------------------------------------------------
\newcommand{\Df}{\mathcal{D}_f}          % forget set
\newcommand{\Dr}{\mathcal{D}_r}          % retain set
\newcommand{\Dra}{\mathcal{D}_r^{\mathrm{adj}}}
\newcommand{\Drr}{\mathcal{D}_r^{\mathrm{rem}}}
\newcommand{\gf}{g_f}

%% ---- result placeholders (Sec. 5 shell) -----------------------------------
\newcommand{\tblph}{\textcolor{gray}{--}}      % empty table cell, result pending
\newcommand{\resph}[1]{\textcolor{gray}{[#1]}} % inline result placeholder
\providecommand{\TailHeadline}{\resph{Tail concentration, semantic coherence,
and the largest observed boundary condition.}}
\providecommand{\PredictionHeadline}{\resph{Prediction claim-pass count,
least-favorable endpoint-gain bound, and principal failure mode.}}
\providecommand{\FidelityHeadline}{\resph{Frozen loss-shake operating point,
rank fidelity, top-$K_p$ overlap, time, memory, and integrity coverage.}}
\providecommand{\ProtectionHeadline}{\resph{Protection claim-pass count,
worst comparator bound, constraint slack, and principal failure mode.}}
\providecommand{\TransferHeadline}{\resph{Predeclared setting coverage,
least-favorable bounds, and settings that do not support the claim.}}
\InputIfFileExists{sections/generated/results_macros.tex}{}{}

\begin{document}

%% Title kept VERBATIM as registered at the abstract deadline (KDD CFP:
%% "Large changes to the title and abstract may result in a desk rejection").
%% The channel framing lives in the abstract's characterization sentence and
%% the body, not in the title.
\title{When Does LLM Unlearning Fail? Predicting and Protecting Susceptible Retained Behavior}

%% Double-blind: no author information in the submission PDF.
\author{Anonymous Author(s)}
\affiliation{%
  \institution{Anonymous Institution}
  \city{}
  \country{}}
\email{anonymous@example.com}

\renewcommand{\shortauthors}{Anon.}

%% ==========================================================================
\begin{abstract}
An LLM asked to forget specific information should not lose knowledge or
capabilities that appear unrelated on the surface. Yet approximate
unlearning can cause such collateral damage. A single average score can
hide this damage when failures concentrate in a small but coherent subset
of retained behaviors. We therefore study whether susceptible retained
behaviors can be identified before unlearning and protected under a fixed
intervention budget. We model susceptibility using two complementary
signals: direction-agnostic local loss sensitivity and hidden-state
proximity to the forget set. For the gradient component, we define
\emph{loss-shake energy} as a pool-shared finite-difference estimator of
classical block-local gradient energy. Unlike
forget-update alignment, loss-shake energy is independent of the direction
ultimately taken by an unlearning trajectory. Paired perturbations estimate
it for an entire candidate pool using batched forward evaluations without
candidate-side backpropagation, while the representation component uses
forward-only hidden-state neighbors. We seal their joint profile before an
unchanged parent unlearner runs, admit only checkpoints passing an external
forgetting gate, and test two independent uses: held-out damage prediction and
fixed-budget repair allocation. Across multiple model families, objectives,
and benchmarks, the evaluation separately requires estimator fidelity,
positive prediction beyond both components and simple controls, lower mean
and worst-tail damage than four allocation controls, and non-inferior native
retained behavior under common forgetting and utility constraints. This
probe-to-protection framework links pre-unlearning risk interpretation to
targeted preservation without changing the parent objective or using realized
damage to select a checkpoint or repair pool.
\end{abstract}


%% ---- CCS / keywords ------------------------------------------------------
\begin{CCSXML}
<ccs2012>
 <concept>
  <concept_id>10002951.10003227.10003351</concept_id>
  <concept_desc>Information systems~Data mining</concept_desc>
  <concept_significance>500</concept_significance>
 </concept>
 <concept>
  <concept_id>10010147.10010257</concept_id>
  <concept_desc>Computing methodologies~Machine learning</concept_desc>
  <concept_significance>500</concept_significance>
 </concept>
 <concept>
  <concept_id>10002978.10002991.10002995</concept_id>
  <concept_desc>Security and privacy~Privacy protections</concept_desc>
  <concept_significance>300</concept_significance>
 </concept>
</ccs2012>
\end{CCSXML}

\ccsdesc[500]{Information systems~Data mining}
\ccsdesc[500]{Computing methodologies~Machine learning}
\ccsdesc[300]{Security and privacy~Privacy protections}

\keywords{machine unlearning, large language models, trustworthy AI,
collateral damage, retained-behavior susceptibility, loss-shake energy,
pre-unlearning risk prediction}

\maketitle

%% ==========================================================================
\section{Introduction}
\label{sec:intro}

A deployed language model may be asked to remove a person's private records,
a copyrighted document, or a harmful capability long after training has
finished. Retraining from scratch for every such request is usually
impractical, so machine unlearning suppresses the designated influence
through targeted post-training updates while preserving the model's
remaining behavior \cite{cao2015towards,bourtoule2021machine,
liu2024rethinking}. The selective requirement is what makes unlearning
hard: forget enough to satisfy the request, but not so aggressively that
useful knowledge is destroyed with it.

The resulting collateral damage is easy to miss because it concentrates.
A model can preserve most retained examples while failing sharply on a
small, coherent subset -- other authors when one author is deleted, benign
science that shares terminology with a suppressed hazardous procedure.
Split-aware evaluations expose this by reporting forget-adjacent retention
separately \cite{cheng2026entanglement,chang2025retain}, but they are
retrospective: the vulnerable behavior is identified only after unlearning
has been run and paid for. Practitioners would rather know \emph{before}
unlearning where damage will land, so that a limited protection or repair
budget can be spent where it is needed.

Prior attempts at such prospective prediction share a hidden assumption:
that susceptibility is a property of the model and data alone, so that some
single score -- lexical or embedding similarity to the forget set, gradient
alignment, gradient magnitude -- should rank vulnerable candidates
regardless of which unlearning method is applied. Empirically these scores
behave inconsistently across methods, and each new proposal appears to win
on some objectives and lose on others. This paper explains that
inconsistency and turns it into structure:

\noindent\emph{Retain-side susceptibility is not a model-intrinsic scalar.
It is conditional on the damage channel induced by the unlearning objective
-- declared in advance by where its loss reads the model.}

\begin{figure*}[t]
\centering
% Figure 1: objective-conditioned damage -> sealed profiles -> crossed test.
% Candidate intensities are deliberately schematic.  The palette is
% grayscale plus one muted blue-gray accent so the logic survives printing.
\definecolor{accent}{HTML}{5E7180}
\definecolor{ink}{HTML}{333333}

\begin{tikzpicture}[
  x=1cm,y=1cm,
  font=\sffamily\scriptsize,
  >={Latex[length=1.65mm,width=1.05mm]},
  paneltitle/.style={font=\sffamily\scriptsize\bfseries,text=ink,anchor=west},
  panelsub/.style={font=\sffamily\scriptsize,text=gray!78!black,anchor=west},
  objective/.style={draw=gray!54,rounded corners=.65mm,fill=white,
                    minimum width=1.58cm,minimum height=1.02cm,line width=.52pt},
  profile/.style={draw=gray!54,rounded corners=.65mm,fill=white,
                  minimum width=1.42cm,minimum height=.98cm,line width=.52pt,
                  align=center},
  candidate/.style={draw=gray!48,rounded corners=.35mm,fill=white,
                    minimum width=.44cm,minimum height=.60cm,inner sep=0pt,
                    line width=.42pt,font=\sffamily\tiny,text=ink},
  seal/.style={draw=accent,rounded corners=.35mm,fill=accent!5,
               inner xsep=.65mm,inner ysep=.45mm,line width=.52pt,
               font=\sffamily\tiny\bfseries,text=accent},
  flow/.style={->,draw=gray!70!black,line width=.62pt},
  trajectory/.style={->,draw=ink!82,line width=.64pt},
  cell/.style={draw=gray!48,fill=gray!3,minimum width=1.00cm,
               minimum height=.63cm,inner sep=0pt,line width=.48pt}
]

\path[use as bounding box] (0,.16) rectangle (17.35,4.72);

% Quiet panel separators.
\draw[gray!25,line width=.4pt] (5.55,.34) -- (5.55,4.58);
\draw[gray!25,line width=.4pt] (11.15,.34) -- (11.15,4.58);

% -------------------------------------------------------------------------
% (a) The same candidates can face different objective-conditioned damage.
\node[paneltitle] at (.16,4.50) {(a)\quad OBJECTIVE-CONDITIONED DAMAGE};
\node[panelsub] at (.16,4.10)
  {same $\theta_0$ and $\mathcal C$; different damage futures};

% Output-reading objective.
\node[objective] (uout) at (.98,3.14) {};
\node[font=\sffamily\tiny\bfseries,text=ink] at (.98,3.48)
  {$\mathcal U_{\mathrm{out}}$};
\foreach \x/\h in {.55/.12,.70/.22,.85/.31,1.00/.18}
  \draw[accent,line width=1.15pt] (\x,2.78) -- (\x,{2.78+\h});
\draw[gray!44,line width=.35pt] (.48,2.78) -- (1.07,2.78);
\node[anchor=west,align=left,font=\sffamily\tiny,text=ink] at (1.12,2.97)
  {loss reads\\[-.35mm]\textbf{outputs}};

% Representation-reading objective.
\node[objective] (urep) at (.98,1.54) {};
\node[font=\sffamily\tiny\bfseries,text=ink] at (.98,1.88)
  {$\mathcal U_{\mathrm{rep}}$};
\foreach \x in {.51,.68,.85}
  \draw[draw=accent,fill=white,rounded corners=.12mm,line width=.4pt]
    (\x,1.18) rectangle ++(.11,.39);
\foreach \y in {1.25,1.37,1.49}
  \fill[accent] (1.02,\y) circle (.75pt);
\node[anchor=west,align=left,font=\sffamily\tiny,text=ink] at (1.13,1.37)
  {loss reads\\[-.35mm]\textbf{hidden states}};

\draw[flow] (uout.east) -- (2.05,3.14);
\draw[flow] (urep.east) -- (2.05,1.54);

% Identical candidate locations; mini-bar height changes the qualitative
% ordering.  The bars carry no numeric experimental value.
\foreach \x/\lab/\height in {2.30/1/.20,2.86/2/.06,3.42/3/.15,3.98/4/.04,4.54/5/.10}{
  \node[candidate] at (\x,3.14) {$x_{\lab}$};
  \fill[gray!68] ({\x-.065},2.85) rectangle ({\x+.065},{2.85+\height});
}
\foreach \x/\lab/\height in {2.30/1/.06,2.86/2/.20,3.42/3/.04,3.98/4/.15,4.54/5/.10}{
  \node[candidate] at (\x,1.54) {$x_{\lab}$};
  \fill[gray!68] ({\x-.065},1.25) rectangle ({\x+.065},{1.25+\height});
}
\node[font=\sffamily\tiny\bfseries,text=ink] at (3.36,.66)
  {same candidates, different damage orderings};

% -------------------------------------------------------------------------
% (b) Both profiles are prospective and use the same candidate universe.
\node[paneltitle] at (5.78,4.50) {(b)\quad PROFILES BEFORE UNLEARNING};
\node[panelsub] at (5.78,4.10)
  {same $\mathcal C$; computed at $\theta_0$ and sealed};

\node[profile] (qg) at (6.55,3.14) {};
\node[font=\sffamily\tiny\bfseries,text=ink] at (6.55,3.48)
  {$\widehat q_G$};
\node[font=\sffamily\tiny,text=ink] at (6.55,3.15)
  {$\ell^-\!\leftarrow\!\theta_0\!\rightarrow\!\ell^+$};
\node[font=\sffamily\tiny,text=ink] at (6.55,2.84)
  {shake energy};

\node[profile] (qh) at (6.55,1.54) {};
\node[font=\sffamily\tiny\bfseries,text=ink] at (6.55,1.88)
  {$\widehat q_H$};
\node[circle,draw=accent,fill=white,minimum size=3.0mm,inner sep=0pt,
      line width=.42pt,font=\sffamily\tiny,text=ink] at (6.55,1.56) {$x$};
\foreach \x/\y in {6.27/1.66,6.82/1.68,6.79/1.39}
  \node[rectangle,draw=accent,fill=accent,minimum size=2.6pt,inner sep=0pt]
    at (\x,\y) {};
\node[font=\sffamily\tiny,text=ink] at (6.55,1.22)
  {hidden neighbors};

\node[font=\sffamily\scriptsize,text=gray!78!black] at (8.82,3.72)
  {shared $\pm$ loss shakes; no candidate backward};
\node[font=\sffamily\scriptsize,text=gray!78!black] at (8.82,2.12)
  {nearest forget-state neighbors};
\draw[flow] (qg.east) -- (7.35,3.14);
\draw[flow] (qh.east) -- (7.35,1.54);

% Accent mini-bars denote prospective exposure, not observed damage.
\foreach \x/\lab/\height in {7.58/1/.18,8.11/2/.07,8.64/3/.15,9.17/4/.05,9.70/5/.10}{
  \node[candidate] at (\x,3.14) {$x_{\lab}$};
  \fill[accent!82] ({\x-.065},2.85) rectangle ({\x+.065},{2.85+\height});
}
\foreach \x/\lab/\height in {7.58/1/.07,8.11/2/.18,8.64/3/.05,9.17/4/.15,9.70/5/.09}{
  \node[candidate] at (\x,1.54) {$x_{\lab}$};
  \fill[accent!82] ({\x-.065},1.25) rectangle ({\x+.065},{1.25+\height});
}
\node[seal] at (10.48,3.14) {SEALED};
\node[seal] at (10.48,1.54) {SEALED};
\node[font=\sffamily\tiny\bfseries,text=ink] at (8.48,.66)
  {two prospective hypotheses, no outcome access};

% -------------------------------------------------------------------------
% (c) Independent trajectories populate a crossed, preregistered endpoint.
\node[paneltitle] at (11.38,4.50) {(c)\quad PREREGISTERED CROSSED TEST};
\node[panelsub] at (11.38,4.10)
  {test, not result: does the relative winner flip?};

% The target trajectories never consume the prospective profiles.
\node[draw=gray!58,rounded corners=.45mm,fill=white,minimum width=.62cm,
      minimum height=.48cm,line width=.5pt,font=\sffamily\tiny\bfseries]
      (theta) at (11.82,2.37) {$\theta_0$};
\node[font=\sffamily\scriptsize,text=gray!78!black] at (12.65,3.72)
  {independent runs};
\node[draw=gray!55,rounded corners=.4mm,fill=white,minimum width=.62cm,
      minimum height=.46cm,line width=.48pt,font=\sffamily\tiny]
      (co) at (12.62,3.12) {$\mathcal U_{\mathrm{out}}$};
\node[draw=gray!55,rounded corners=.4mm,fill=white,minimum width=.62cm,
      minimum height=.46cm,line width=.48pt,font=\sffamily\tiny]
      (cr) at (12.62,1.62) {$\mathcal U_{\mathrm{rep}}$};
\node[draw=gray!55,rounded corners=.4mm,fill=gray!4,minimum width=.58cm,
      minimum height=.46cm,line width=.48pt,font=\sffamily\tiny]
      (do) at (13.57,3.12) {$d_{\mathrm{out}}$};
\node[draw=gray!55,rounded corners=.4mm,fill=gray!4,minimum width=.58cm,
      minimum height=.46cm,line width=.48pt,font=\sffamily\tiny]
      (dr) at (13.57,1.62) {$d_{\mathrm{rep}}$};
\draw[trajectory] (theta.east) -- (12.14,2.37) |- (co.west);
\draw[trajectory] (theta.east) -- (12.14,2.37) |- (cr.west);
\draw[trajectory] (co) -- (do);
\draw[trajectory] (cr) -- (dr);

% Correlation matrix.  All four cells remain visually neutral because the
% preregistered interaction is a question, not an observed result.
\node[font=\sffamily\scriptsize,text=gray!78!black,align=center] at (15.23,3.57)
  {output\\damage};
\node[font=\sffamily\scriptsize,text=gray!78!black,align=center] at (16.39,3.57)
  {representation\\damage};
\node[font=\sffamily\tiny\bfseries,text=ink,align=right] at (14.42,2.82)
  {$\widehat q_G$};
\node[font=\sffamily\tiny\bfseries,text=ink,align=right] at (14.42,2.02)
  {$\widehat q_H$};
\node[cell]        at (15.23,2.82) {$\rho_{G,\mathrm{out}}$};
\node[cell]        at (16.39,2.82) {$\rho_{G,\mathrm{rep}}$};
\node[cell]        at (15.23,2.02) {$\rho_{H,\mathrm{out}}$};
\node[cell]        at (16.39,2.02) {$\rho_{H,\mathrm{rep}}$};

\node[font=\sffamily\scriptsize,text=ink] at (15.17,.79)
  {$\displaystyle
    \Delta=
    (\rho_{G,\mathrm{out}}-\rho_{H,\mathrm{out}})
    -(\rho_{G,\mathrm{rep}}-\rho_{H,\mathrm{rep}})>0\;?$};

\end{tikzpicture}

\caption{\textbf{Susceptibility is tested as an objective-conditioned
interaction.} \textbf{(a)} For the objectives studied, the direct
forget-suppression term reads either outputs or hidden states; the same
$\theta_0$ and retained candidates may therefore yield different damage
rankings. \textbf{(b)} The loss-shake gradient-energy and hidden-state-proximity
hypotheses score the same candidates at $\theta_0$ and are sealed before any
outcome is observed. \textbf{(c)} Independent unlearning trajectories that
never access the scores produce the two damage columns; matrix entries are
Spearman correlations. The preregistered endpoint asks whether the relative
predictive advantage crosses over between channels, $\Delta>0$, not whether
one probe is universally superior. Candidate intensities are schematic, not
experimental results.}
\Description{A three-panel schematic with no observed results. Panel (a)
applies an output-reading and a hidden-state-reading unlearning objective to
the same initial model and retained-candidate set; qualitative bands under
the repeated candidate identifiers illustrate two possible damage orderings.
Panel (b) applies loss-shake gradient-energy and hidden-state-proximity
profiles to that same candidate set and marks both profiles as sealed before
unlearning. Panel (c) shows two independent objective trajectories producing
output-channel and representation-channel damage, followed by a two-by-two
matrix of correlations between each sealed profile and each damage outcome.
The difference-in-differences expression below the matrix asks whether the
relative predictive winner crosses over between channels.}
\label{fig:channel-mechanism}
\end{figure*}

To first order, the damage to a retained candidate $x$ under a parameter
update $\Delta\theta$ is $d(x)\approx g_x^{\top}\Delta\theta$, which makes
$\lVert g_x\rVert$ a direction-free exposure envelope rather than a
guarantee of realized damage. What
$\Delta\theta$ looks like is determined by the objective's loss. Objectives
whose loss is a function of model \emph{outputs} -- token likelihoods or
preferences, as in GA, GradDiff, NPO, SimNPO, IdkDPO, and GRU
\cite{yao2024large,zhang2024negative,fan2024simplicity,maini2024tofu,
wang2025gru} -- motivate gradient-energy exposure. Objectives whose loss is a
function of internal \emph{representations} -- hidden states, as in RMU,
RepNoise, and Circuit Breakers
\cite{li2024wmdp,rosati2024repnoise,zou2024circuit} -- displace activation
geometry and motivate proximity to forget representations. The two
channels therefore call for different predictor families. We test that
claim at the family level with a frozen objective-by-probe interaction;
individual objectives may show partial, mixed, or reversed signal.

Acting on this structure is only useful if the matching signal is cheap at
candidate-pool scale. Exact per-candidate gradient norms require one
backward pass per candidate. We instead estimate the entire pool's gradient
magnitudes with a \emph{backward-free randomized-sensitivity probe}:
symmetric forward evaluations at $\theta_0\pm\eta v_r$ along a few shared
random directions $v_r$ turn squared central differences into an unbiased
(up to $O(\eta^2)$) estimate of every candidate's gradient energy
simultaneously, with Monte Carlo variance controlled by the number of
shared directions. A
fidelity gate separates the estimator's numerical validity from its
predictive value. The representation channel uses a complementary
forward-only hidden-state proximity profile.

Prediction becomes protection through routing. Given a deletion request and
a declared objective, the channel determines the probe; the probe ranks the
discovery-fold candidates; the top-ranked pool receives a fixed repair
budget under a guarded procedure that caps forgetting regression. Our
success criterion is deliberately not a leaderboard: it is the
\emph{interaction} in a preregistered crossed design. For an output-channel
parent (GradDiff) and a representation-channel parent (RMU), we cross each
with matched, mismatched, random, and no-protection selectors and test
whether the channel-matched selector yields the lowest audit-tail damage at
matched forgetting -- per parent. All arms and reach failures are reported;
the channel claim is not declared successful unless both parent-wise tests
pass.

Our contributions are threefold:
\begin{itemize}
    \item \textbf{Channel-conditioned susceptibility.} We test whether the
    predictive signal changes with the objective's channel, declared from
    its loss definition. Evidence comes from a frozen family-level
    interaction over a channel-stratified objective roster, not a
    per-method anecdote.

    \item \textbf{A backward-free probe for the gradient channel.} A
    randomized symmetric-perturbation estimator profiles gradient
    magnitude over the whole candidate pool using only $2R$ batched
    forward sweeps, with a preregistered fidelity gate that certifies when
    the estimate is trustworthy under a given radius and precision.

    \item \textbf{A falsifiable protection test.} A crossed design routes
    the same repair budget through matched, mismatched, random, and no-score
    selectors, testing whether the pre-unlearning channel distinction is
    operationally useful at matched forgetting.
\end{itemize}

We position this work as a prediction-and-protection \emph{layer} on top of
existing unlearners, not as a competing unlearning objective. The remainder
of the paper follows this structure. Section~\ref{sec:related} organizes
prior unlearning methods by channel and identifies the channel-matching
gap. Section~\ref{sec:setup} formalizes damage channels, the prediction
and protection endpoints, and the preregistered interaction criterion
(Figure~\ref{fig:channel-mechanism}). Section~\ref{sec:method} presents
the probes, the fidelity gate, and the guarded channel-matched protection
procedure. Section~\ref{sec:experiments} reports the channel matrix, the
crossed protection experiment, cost, and cross-dataset replication.

\section{Related Work}
\label{sec:related}

\paragraph{Output-channel unlearning: losses on model outputs.}
One family of LLM unlearning objectives defines its loss on model
\emph{outputs} -- token likelihoods or preferences over completions.
Gradient ascent and gradient difference directly raise forget-answer NLL
while optionally re-descending on retain data
\cite{jang2023knowledge,yao2024large}. Negative preference optimization and
its simplified variant reformulate forgetting as preference optimization
against the forget answer \cite{zhang2024negative,fan2024simplicity};
IdkDPO substitutes refusal targets \cite{maini2024tofu}; weighted and
saturation-aware variants reshape the same likelihood pressure
\cite{wang2025gradient,yang2025satimp}. Retention-aware members of this
family operate on the same channel: GRU rectifies conflicts between forget
and aggregate-retain \emph{loss gradients} \cite{wang2025gru}, and gradient
surgery variants synthesize the update from the same output-loss geometry
\cite{zhou2025implicit,lin2024gdr,patel2025learning,wu2025munba,
xiao2026sago}. Whatever their differences, all of these methods move
parameters along directions assembled from output-loss gradients.

\paragraph{Representation-channel unlearning: losses on hidden states.}
A second family defines its loss on internal \emph{representations}. RMU
steers forget-set hidden states toward a random control vector while
pinning retain-set hidden states \cite{li2024wmdp}. RepNoise pushes
harmful-content representations toward noise across layers
\cite{rosati2024repnoise}, and Circuit Breakers reroute activations that
enter harmful regions toward orthogonal directions \cite{zou2024circuit}.
Localized-editing approaches that overwrite specific weight or neuron
subspaces \cite{meng2022locating,fan2024salun,foster2024ssd,cai2025slug}
share this character: the optimized quantity is internal structure, and
output change is a consequence rather than the loss argument. Benchmarks --
TOFU, MUSE, WMDP, RWKU, KnowUnDo, OpenUnlearning
\cite{maini2024tofu,shi2024muse,li2024wmdp,jin2024rwku,tian2024forget,
dorna2025openunlearning} -- evaluate members of both families, and
robustness studies show that endpoint forget metrics need not certify
removal \cite{lynch2024eight,patil2024can,hu2024jogging,
lucki2024adversarial,deeb2024unlearning}.

\paragraph{The channel-matching gap.}
Retain-side risk has so far been diagnosed \emph{within} one family's own
geometry. G-effect analyzes output-loss gradient interactions across steps
and layers \cite{wang2025gradient}; GUARD scores forget examples by
alignment with a retain Jacobian \cite{niu2025guard}; HAMU tracks an
aggregate forget--retain gradient dot product \cite{chen2026hamu};
split-aware evaluation and constrained recovery protect a forget-adjacent
retain slice identified by proximity \cite{cheng2026entanglement}. Each of
these commits to a single signal family -- gradient geometry or similarity
-- independent of which objective will actually be run. To our knowledge,
no prior work asks whether the \emph{right retain-side predictor depends on
the objective's damage channel}, tests that dependence with balanced
objective rosters on both channels, or routes protection through the
matching signal. That interaction is precisely where single-score methods
have been reported to behave inconsistently, and it is the gap this paper
fills: a declared, loss-definition-level classification of objectives into
channels, matched predictor families, and a crossed protection experiment
that tests whether matching is operationally necessary.

\paragraph{Attribution and randomized probing.}
Influence functions, TracIn, TRAK, LESS, and DataInf estimate
training-data influence but require per-example gradients or curvature
machinery \cite{koh2017influence,pruthi2020tracin,park2023trak,
xia2024less,kwon2024datainf}; Forward-INF and For-Value amortize the
expensive side over candidates with forward passes
\cite{ko2024mirrored,deng2026forvalue}, though their estimand is
retrospective attribution rather than prospective unlearning risk. On the
numerical side, SPSA, random gradient-free methods, MeZO, and randomized
trace estimators establish that random symmetric perturbations recover
gradient information from forward evaluations alone
\cite{spall1992spsa,nesterov2017random,malladi2023mezo,
hutchinson1989stochastic}. We repurpose this machinery, not for
optimization, but to score every retained candidate's
$\lVert g_x\rVert_2$ simultaneously from $2R$ shared forward sweeps --
giving the output channel a predictor whose cost matches the forward-only
proximity signal available to the representation channel. The estimator is
classical; the contributions are its role as a prospective, channel-matched
unlearning-risk profile, the fidelity gate that certifies it, and the
demonstration that acting on it protects the audit tail.

\section{Problem Setup and Damage Channels}
\label{sec:setup}

\subsection{Deletion Request and Retained-Behavior Scope}

Let $\theta_0$ denote the model before unlearning. A deletion request
specifies a forget set $\mathcal D_f$ and an enumerated universe
$\mathcal C$ of retained behaviors that should remain usable. An element
$x\in\mathcal C$ is a prompt--answer behavior on which retention can be
measured; it need not be a record from the original training set.

We split $\mathcal C$ at the semantic-group level into disjoint discovery
and audit folds,
$\mathcal C=\mathcal C_{\mathrm{disc}}\,\dot\cup\,\mathcal C_{\mathrm{audit}}$.
All profile-guided selection and protection use only
$\mathcal C_{\mathrm{disc}}$. Scores on $\mathcal C_{\mathrm{audit}}$ may be
computed before unlearning for prediction evaluation, but they are sealed
and cannot affect optimization, hyperparameter selection, checkpoint
selection, or stopping. Where a benchmark provides a native retain
partition, it is included in the sealed audit fold and denoted
$\mathcal A_{\mathrm{native}}$.

For a behavior $x$ and parameters $\theta$, let
\[
    \ell(x;\theta)
    =
    -\frac{1}{|y_x|}
    \sum_{j=1}^{|y_x|}
    \log p_\theta
    \bigl(y_{x,j}\mid p_x,y_{x,<j}\bigr)
\]
be the mean answer-token negative log-likelihood, with prompt tokens
masked. The realized damage to $x$ at a checkpoint $\theta_t$ of an
unlearning procedure is
\begin{equation}
    d_t(x)=\ell(x;\theta_t)-\ell(x;\theta_0),
    \label{eq:realized-damage}
\end{equation}
with positive values indicating harmful retention loss.

\subsection{Damage Channels}
\label{sec:setup-channels}

To first order in the update,
\[
    d(x)\approx g_x^{\top}\Delta\theta,
    \qquad
    |d(x)|\lesssim\lVert g_x\rVert_2\lVert\Delta\theta\rVert_2,
    \qquad
    g_x=\nabla_\theta\ell(x;\theta_0).
\]
This first-order coupling between a candidate's local gradient and the
realized update is the \emph{update-conditioned gradient interference}
mechanism named in the abstract. Gradient magnitude is then a
direction-free \emph{exposure envelope}; it
does not imply damage to every high-score candidate. Whether the envelope
ranks realized damage is an empirical interaction tested across objectives.
The method instantiates the envelope on a declared parameter block $B$
covering the parents' trainable parameters (Section~\ref{sec:method}).
The update is nevertheless structured: its gradients are set by where the
unlearning loss directly reads the model. Current objectives place their
forget-suppression signal at one of two readouts.

\paragraph{Output channel (loss-gradient).}
If the forget-suppression loss directly supervises outputs -- answer-token
likelihoods or preferences, as in GA, GradDiff, NPO, SimNPO, IdkDPO, and
GRU -- then $\Delta\theta$ is assembled from output-loss gradients. Trying
to predict each candidate's signed exposure $g_x^{\top}\Delta\theta$
requires anticipating the realized update direction, which is brittle over
a multi-step trajectory because that direction rotates with the model and
optimizer state. We instead measure the direction-free exposure envelope
$\lVert g_x\rVert_2$: how strongly the candidate loss can respond to small
parameter motion, before knowing which output-channel trajectory will run.

\paragraph{Representation channel.}
If the forget-suppression loss directly supervises internal hidden states,
as in RMU, RepNoise, and Circuit Breakers, the update is chosen to displace
activation geometry in a region occupied by the forget set. The matched
hypothesis is therefore \emph{proximity in hidden-state space}: candidates
whose representations share that region are more exposed to its
displacement. This is complementary to, rather than implied by, an
output-loss gradient norm.

\paragraph{Declaration, not post-hoc labeling.}
The channel of an objective is read off its loss definition before any
experiment is run -- a static inspection of which tensors enter the loss,
requiring no training run: token likelihoods or preferences imply the
output channel; hidden states imply the representation channel. For
objectives that mix terms (RepNoise adds an auxiliary output-side ascent
term), the declaration rule is the \emph{forget-suppression term}: the
loss component that drives removal decides the channel. Likewise, GRU
modifies the update rule but its losses read outputs, so it is declared
output-channel. We preregister the assignment (output: GA, GradDiff,
NPO, SimNPO, IdkDPO, GRU; representation: RMU, RepNoise, Circuit Breakers)
and never revise it against results. Two honesty clauses apply. First, the
dichotomy reflects where \emph{existing} objectives read the model, not a
law of nature; an objective reading some third quantity would define a
third channel, and the framework extends to it. Second, the channels are
poles of a spectrum rather than a partition of downstream effects: an
output-channel objective also perturbs representations incidentally, while
a representation objective ultimately changes outputs. The declaration is
about the objective's direct supervision signal, not every tensor in its
computation graph. Cross-channel signal is therefore measured rather than
ruled out; the confirmatory claim concerns the interaction between declared
channel and probe family.

\subsection{Prospective Prediction and the Interaction Endpoint}
\label{sec:setup-prediction}

A pre-unlearning predictor assigns each candidate a scalar score
$r(x;\theta_0,\mathcal D_f)$ before any target trajectory is run.
Predictive validity is evaluated against damage generated by independently
configured unlearning procedures that receive no access to $r$: within each
objective, we report Spearman correlation between the sealed score and
realized damage, plus tail-focused summaries.

The channel claim is not that some probe achieves high correlation -- it is
that \emph{which} probe family predicts flips with the objective's channel.
The preregistered primary endpoint is therefore a difference-in-differences
interaction. For a gradient-family probe $r_g$, a representation-family
probe $r_h$, an output-channel objective $\mathcal U_{\mathrm{out}}$, and a
representation-channel objective $\mathcal U_{\mathrm{rep}}$, with
$\rho(r,\mathcal U)$ the Spearman correlation of probe $r$ on damage from
$\mathcal U$:
\begin{equation}
    \Delta
    =
    \bigl[\rho(r_g,\mathcal U_{\mathrm{out}})-\rho(r_h,\mathcal U_{\mathrm{out}})\bigr]
    -
    \bigl[\rho(r_g,\mathcal U_{\mathrm{rep}})-\rho(r_h,\mathcal U_{\mathrm{rep}})\bigr].
    \label{eq:interaction}
\end{equation}
$\Delta>0$ with a candidate-bootstrap confidence interval excluding zero
supports channel matching. The double difference cancels both nuisance
factors that contaminate single comparisons: a probe that is uniformly
better, and an objective that is uniformly more predictable.

\subsection{Channel-Matched Protection and the Crossed Design}
\label{sec:setup-protection}

A score profile allocates a limited protection budget: using only the
discovery fold, a selector constructs a protected pool
$\mathcal P_{K_p}(r)\subseteq\mathcal C_{\mathrm{disc}}$ of the $K_p$
highest-ranked eligible behaviors, plus a matched low-score reference pool
$\mathcal R_0(r)$ (construction details in Section~\ref{sec:method}).
Protection methods are compared only at a common forgetting criterion
$\Gamma_f(\theta)\leq\tau_f$; a run that does not reach the criterion is a
reach failure, not a comparison point. Conditional on reach, the
confirmatory outcomes on the sealed audit are mean damage and upper-tail
damage $\operatorname{CVaR}_{.95}(\{d(x)\})$ -- the tail because that is
where collateral damage concentrates and where aggregate metrics hide it.

The operational endpoint is again an interaction, realized as a
preregistered crossed design: each parent objective (one per channel) is
run with each selector (matched family, mismatched family, random, none),
holding the parent, budgets, and repair procedure fixed. Success means the
channel-matched selector attains the lowest audit CVaR per parent. This
criterion deliberately does not ask our pipeline to outperform any
particular unlearner; it asks whether channel information -- and only
channel information -- moves protection where it helps.

% 4.method.tex
% When Does LLM Unlearning Fail?
% Channel-conditioned susceptibility: probes, routing, guarded protection.

\section{Method}
\label{sec:method}

Section~\ref{sec:setup-channels} established which candidate property
couples to each damage channel. This section supplies the operational
pieces: a backward-free estimator of gradient magnitude for the output
channel (Section~\ref{sec:method-gradient}), a forward-only proximity
score for the representation channel
(Section~\ref{sec:method-representation}), the routing and pool
construction that turn a declared objective into a protect pool
(Section~\ref{sec:method-routing}), and the guarded repair procedure that
spends the protection budget without undoing forgetting
(Section~\ref{sec:method-protection}).

\subsection{Gradient-Channel Probe: Randomized Sensitivity}
\label{sec:method-gradient}

Let $B$ be a declared parameter block and
$g_x=\nabla_B\ell(x;\theta_0)$ the candidate gradient at the
pre-unlearning model. The output-channel target quantity is the
direction-agnostic sensitivity
\begin{equation}
    m(x)=\lVert g_x\rVert_2.
    \label{eq:exact-susceptibility}
\end{equation}
Computing $m(x)$ exactly requires one backward evaluation per candidate;
we retain this exact profile as the expensive reference. The scalable
estimator reuses symmetric forward perturbations. Let
$v_1,\ldots,v_R$ be seeded random unit vectors in $B$, and let
\begin{equation}
    \delta_\eta(x;v)
    =
    \frac{\ell(x;\theta_0+\eta v)-\ell(x;\theta_0-\eta v)}{2\eta}
    \label{eq:directional-probe}
\end{equation}
be the symmetric response of candidate $x$ along $v$. The randomized
sensitivity score is
\begin{equation}
    r_{\mathrm{rand}}(x)
    =
    \frac{1}{R}
    \sum_{r=1}^{R}
    \delta_\eta(x;v_r)^2 .
    \label{eq:randomized-susceptibility}
\end{equation}
Because $\mathbb E_v[(g_x^\top v)^2]=\lVert g_x\rVert_2^2/d_B$ for random
unit directions in a $d_B$-dimensional block, and the central difference
cancels the leading curvature term,
$r_{\mathrm{rand}}(x)\approx\lVert g_x\rVert_2^2/d_B$ up to Monte Carlo
and $O(\eta^2)$ error; the shared factor $d_B$ does not affect ranking.
The relative Monte Carlo variance is $2/R$ independently of $d_B$, so a
few dozen directions suffice. The directions are shared across the entire
candidate universe: $R$ directions cost $2R$ batched candidate-pool
forward sweeps in total, versus one backward evaluation \emph{per
candidate} for the exact profile. Memory never exceeds inference: no
per-example gradient is materialized.

\paragraph{Fidelity gate.}
The estimator is trusted only where a preregistered gate certifies it. The
radius $\eta$ must sit between two failure regimes: too small, and the
paired losses agree to within floating-point resolution so their
difference is cancellation noise; too large, and higher-order terms
contaminate the response. The gate first verifies the perturbation is
realized at all (measured $\lVert\eta v\rVert/\eta$ and the fraction of
changed block parameters, which catches low-precision underflow), then
decomposes agreement between three quantities on a diagnostic subset: the
exact squared norm $A=\lVert g_x\rVert^2$, the exact projected proxy
$B=(d_B/R)\sum_r (g_x^\top v_r)^2$, and the finite-difference estimate
$C$. Rank agreement of $B$ with $A$ isolates Monte Carlo error; rank
agreement of $C$ with $B$ isolates finite-difference error, including
cancellation. The radius and $R$ are frozen once the gate passes and
before any target trajectory is unsealed.

\paragraph{Rejected diagnostic: forget-direction alignment.}
The same probe evaluated along the normalized initial forget gradient
$\hat g_f$ yields the classical alignment score
$\delta_\eta(x;\hat g_f)\approx\langle g_x,\hat g_f\rangle$. We verified
this estimate against exact Jacobian--vector products, confirmed its
numerical validity, and preregistered its rejection as a susceptibility
predictor: across engines it correlates weakly or negatively with realized
damage, because a multi-step trajectory does not preserve its initial
direction. It appears in our tables as a cautionary baseline; the
supporting analysis is in the appendix.

\subsection{Representation-Channel Probe: Hidden-State Proximity}
\label{sec:method-representation}

The representation channel damages candidates whose hidden states occupy
the displaced region, so its matched probe is proximity, not gradient
magnitude. Candidate and forget examples are embedded with mean-pooled
final hidden states of the pre-unlearning model; each candidate is scored
by its mean cosine similarity to its $k$ nearest forget examples. The
score is forward-only, needs no parameter perturbation, and -- decisively
for the crossed design -- selects an almost disjoint candidate set from
the gradient probe: the two families' rankings are nearly uncorrelated on
our benchmarks. External sentence embeddings and lexical overlap provide
weaker members of the same family and are reported as comparators.

\subsection{Channel Routing and Profile Construction}
\label{sec:method-routing}

Given a deletion request and a declared parent objective, the router reads
the objective's channel from its loss definition
(Section~\ref{sec:setup-channels}) and selects the matched probe family:
randomized sensitivity for output-channel parents, hidden-state proximity
for representation-channel parents. Nothing in the router depends on
outcomes; it is a lookup on the preregistered channel table.

The selected scorer assigns one value to every candidate in the declared
universe. For a fixed budget $K_p$, the protect pool is
$\mathcal P_{K_p}(r)=\operatorname{TopK}_{x\in\mathcal C_{\mathrm{disc}}}\,r(x)$,
with deterministic tie-breaking; a size-matched low-score pool
$\mathcal R_0(r)$ is sampled from a frozen near-zero quantile of the
discovery fold and serves as the repair-time reference stream. Pool size,
quantile, group-disjoint folds, seed, and fallback rules are fixed before
target outcomes are observed, and no audit candidate can enter either
pool.

\subsection{Soft Channel Routing: The Channel-Mixture Score}
\label{sec:method-mixture}

The discrete router commits fully to one probe family. The spectrum
clause of Section~\ref{sec:setup-channels}, however, predicts residual
off-channel signal -- output-channel objectives perturb representations
incidentally -- so we generalize routing to a one-parameter family. Let
$\tilde r_g,\tilde r_h:\mathcal C\to[0,1]$ be the two headline probes
after rank normalization over the candidate universe (rank transforms
because the raw scales -- a squared-gradient estimate and a cosine -- are
not commensurable; the normalization is part of the frozen definition).
The \emph{channel-mixture score} is
\begin{equation}
    s_\alpha(x)
    =
    (1-\alpha)\,\tilde r_g(x)
    +
    \alpha\,\tilde r_h(x),
    \qquad \alpha\in[0,1],
    \label{eq:mixture}
\end{equation}
so that $\alpha$ is the weight on the representation route. The
endpoints recover the discrete router exactly -- $\alpha{=}0$ is the
randomized-sensitivity probe, $\alpha{=}1$ the proximity probe -- so
every result of the preceding sections is the $\alpha\in\{0,1\}$ special
case, and the declared channel now supplies a \emph{prior} over $\alpha$
rather than a hard switch: output-channel parents map to $\alpha{=}0$,
representation-channel parents to $\alpha{=}1$, and hybrid objectives
(e.g.\ RepNoise's auxiliary output term) to interior values.

Two disciplines keep the mixture honest. First, $\alpha$ is never chosen
on sealed data: it is either (i) the preregistered channel-to-$\alpha$
mapping, fixed before any outcome is read (primary -- the deployed
router is unchanged, and the mixture is preregistered as exploratory
instrumentation that does not reweight the primary endpoints), or
(ii) fit on discovery folds only (adaptive variant), with the sealed
audit untouched in both cases. Second, prediction and protection are allowed different
optima: the crossed experiment of Section~\ref{sec:exp-protection}
motivates measuring the prediction-optimal $\alpha^{*}_{\mathrm{pred}}$
(which maximizes damage correlation) separately from the
protection-optimal $\alpha^{*}_{\mathrm{prot}}$ (which minimizes sealed
audit tail damage when the pool is $\operatorname{TopK}$ by $s_\alpha$),
because repair generalization can propagate through representation-space
neighborhoods regardless of the parent's channel. The mixture sweep of
Section~\ref{sec:exp-mixture} reports both curves; how far the two
optima separate is itself an empirical finding, not an assumption.

\subsection{Guarded Channel-Matched Repair}
\label{sec:method-protection}

The budget is spent by an engine-then-repair pipeline that leaves the
parent unlearner untouched. The parent runs from $\theta_0$ until its
first saved checkpoint satisfying the common forget criterion; a parent
that never reaches remains a reach failure and is not repaired. At the
reaching state we cache per-sequence and per-token forget losses as paired
references, then minimize answer-token NLL on $\mathcal P_{K_p}$.

Repair must not silently undo forgetting. At fixed refresh intervals we
compute mean forget and reference-stream gradients and project the repair
update away from their span, making those mean losses locally invariant to
first order. Because a mean-level projection does not protect every forget
target, an identity-paired one-sided guard supplements it: with
$\bar u_i$ the cached post-unlearning forget loss and $u_i$ its current
value,
\begin{equation}
    D_{\downarrow}(u,\bar u)
    =
    \frac{1}{n}
    \sum_{i=1}^{n}
    \bigl[\bar u_i-\epsilon-u_i\bigr]_+^2
    \label{eq:guard}
\end{equation}
penalizes any per-example forgetting regression beyond tolerance
$\epsilon$, at sequence and (where enabled) token level. At audited
refreshes, a repair step that exceeds either paired-loss budget is
rejected, the last accepted parameters are restored, and the step size is
reduced. The output is the final saved repair checkpoint that still
satisfies the forget criterion; retention is evaluated only on the sealed
audit.

In the crossed design of Section~\ref{sec:setup-protection}, every
selector arm -- matched, mismatched, random -- runs this identical
pipeline with the identical budget; only the ranking that fills
$\mathcal P_{K_p}$ changes. Any difference between arms is therefore
attributable to where the budget went, which is exactly the channel
question.

\subsection{Scope of the Method}
\label{sec:method-scope}

The profile is a prospective risk measure, not a deletion certificate:
forgetting efficacy must still be established by direct, paraphrase, and
adversarial audits. The channel table covers where current objectives read
the model; an objective reading a third quantity would require a third
probe family, and hybrid objectives inherit partial signal from both. The
backward-free claim is scoped to candidate-side profiling: parents and
repair still backpropagate, as they must.

% ================================================================= Experiments

\section{Experiments}
\label{sec:experiments}

The experiments instantiate the two preregistered interaction endpoints of
Section~\ref{sec:setup} and the fidelity claims of
Section~\ref{sec:method}. All endpoints, channel assignments, probe
settings ($\eta$, $R$), and the dataset roster were frozen before the
confirmatory runs were unsealed. The evaluation addresses four questions:

\begin{enumerate}
    \item \textbf{Channel-conditioned prediction:} Across a balanced
    roster of output-channel and representation-channel objectives, does
    the winning predictor family flip with the channel
    (interaction $\Delta>0$, Eq.~\eqref{eq:interaction})?

    \item \textbf{Backward-free fidelity:} How much of the exact
    gradient-norm profile does randomized sensitivity recover, at what
    cost, and where does the fidelity gate draw its validity boundary?

    \item \textbf{Channel-matched protection:} In the crossed design
    (parents $\times$ selectors), does the matched selector minimize
    sealed-audit tail damage per parent, at matched forgetting?

    \item \textbf{Robustness:} Does the interaction survive across
    datasets whose forget--retain proximity structure differs, across
    seeds and requests, and under the declared honesty clauses?
\end{enumerate}

Table~\ref{tab:channel-matrix} is the confirmatory prediction experiment,
Table~\ref{tab:crossed} the confirmatory protection experiment, and
Table~\ref{tab:bwfree} the fidelity-and-cost account of the randomized
probe. The strong-method comparison, numerical-fidelity detail, mechanism,
specificity, and complete transfer results appear in the appendix.


\subsection{Experimental Setup}
\label{sec:benchmark-protocol}

\paragraph{Objectives (balanced by channel).}
Output channel: GA, GradDiff, NPO, SimNPO, IdkDPO, GRU. Representation
channel: RMU, RepNoise, Circuit Breakers. All nine parents are configured
on development requests without access to susceptibility scores; their
objectives, trainable blocks, retain streams, learning rates, budgets, and
checkpoint grids are fixed before scoring. Damage is
Eq.~\eqref{eq:realized-damage} at the predeclared horizon, including
requests that fail the forget criterion.

\paragraph{Predictors.}
Gradient family: exact gradient norm (per-candidate backward; expensive
reference) and randomized sensitivity
(Eq.~\eqref{eq:randomized-susceptibility}; $2R$ forward sweeps, ours).
Representation family: hidden-state proximity (matched probe),
sentence-embedding proximity, lexical overlap. Cautionary baselines:
forget-direction alignment (numerically validated, preregistered as
rejected), one random direction, initial candidate NLL, answer length,
random ranking. Every predictor is frozen at $\theta_0$ and sealed before
any trajectory runs.

\paragraph{Datasets and models.}
The primary setting is TOFU single-author requests \cite{maini2024tofu}
with group-disjoint discovery/audit folds and a sealed native audit.
Replication settings re-run the reduced crossed roster
(\{GradDiff, NPO\}$\times$\{RMU, RepNoise\}) on WMDP-bio with an MMLU
retain universe \cite{li2024wmdp,hendrycks2021measuring} (the
representation family's native habitat), MUSE-News \cite{shi2024muse}
(real-world corpus, output-channel habitat), and RWKU \cite{jin2024rwku}
(pretrained knowledge with built-in neighbor audits). Stress settings
probe the honesty clauses: MUSE-Books (verbatim memorization) and PISTOL
\cite{qiu2024pistol} (maximal forget--retain structural entanglement,
where proximity probes are most at risk of over-selecting). Because each
dataset fixes the probes' base rates while the channel is fixed by the
objective, replication across these settings tests that the interaction --
not any single correlation level -- is the invariant. Mechanism runs use
Qwen2.5-1.5B in fp32, where the fidelity gate certifies the randomized
probe; Qwen2.5-7B in bf16 provides the cost table and scale transfer with
the exact-norm probe standing in where low precision fails the gate.

\paragraph{Protocol.}
Prediction metrics are within-request Spearman $\rho$, AUROC against the
realized top-$K_p$ damage set, and tail-restricted $\rho$; the primary
endpoint is the interaction of Eq.~\eqref{eq:interaction} with
candidate-level bootstrap intervals. Protection follows the crossed design
at a common preregistered forget-recall criterion with paraphrase audits;
outcomes are mean and $\operatorname{CVaR}_{.95}$ damage on the sealed
native audit, reach reported unconditionally. Uncertainty uses
hierarchical bootstrap over requests and seeds; leave-one-out ranges
accompany small top levels.

\begin{figure*}[t]
\centering
\begin{minipage}[t]{0.318\textwidth}
\centering
\fbox{\parbox{0.92\linewidth}{\centering\vspace{1.05cm}
\textbf{A. Channel matrix}\\[1mm]
\resph{Heatmap of $\rho$(predictor, objective): rows = predictor families,
columns = objectives grouped by declared channel. Gradient rows light up
left and go dark on RMU; proximity rows peak right.}
\vspace{1.05cm}}}
\end{minipage}\hfill
\begin{minipage}[t]{0.318\textwidth}
\centering
\fbox{\parbox{0.92\linewidth}{\centering\vspace{1.05cm}
\textbf{B. Crossover}\\[1mm]
\resph{Mean $\rho$ of the two headline probes by objective channel: the
lines cross -- the winning probe flips with the channel.}
\vspace{1.05cm}}}
\end{minipage}\hfill
\begin{minipage}[t]{0.318\textwidth}
\centering
\fbox{\parbox{0.92\linewidth}{\centering\vspace{1.05cm}
\textbf{C. Crossed protection}\\[1mm]
\resph{Sealed-audit CVaR per parent $\times$ selector: matched lowest for
both parents; mismatched $\approx$ random $\approx$ none.}
\vspace{1.05cm}}}
\end{minipage}
\caption{\textbf{Main result: susceptibility is channel-conditional and
channel matching is actionable.} \textbf{A:} the objective$\times$family
matrix behind Table~\ref{tab:channel-matrix}. \textbf{B:} the interaction
visualized -- no probe wins on both channels. \textbf{C:} the crossed
protection experiment of Table~\ref{tab:crossed}: routing the same budget
through the channel-matched probe reduces worst-tail audit damage where
mismatched, random, and no protection do not.}
\Description{Three panels: a predictor-by-objective correlation heatmap
grouped by channel, a two-line crossover plot of headline probes across
channels, and grouped bars of audit tail damage for parent-selector
combinations.}
\label{fig:channel-main}
\end{figure*}


\subsection{Main Result I: Channel-Conditioned Prediction}
\label{sec:exp-prediction}

Table~\ref{tab:channel-matrix} reports Spearman $\rho$ between each sealed
predictor and realized damage for every objective, with objectives grouped
by declared channel and predictors grouped by family. The reading is
columnwise-within-rowgroup: gradient-family probes should dominate the
output-channel columns and lose their signal on the representation
columns, with the representation family showing the mirrored pattern; the
per-family channel means and the preregistered interaction $\Delta$ (with
bootstrap CI, reported in the caption once unsealed) quantify the flip.
The asymmetry clause of Section~\ref{sec:setup-channels} predicts partial
representation-probe signal on output-channel columns but near-zero
gradient-probe signal on RMU; the cautionary alignment baseline should
predict nowhere. Flagged columns are retained rather than dropped:
$^{\dag}$ columns measure a weak-signal (near-static) regime and
$^{\ddag}$ marks criterion reach achieved only through model-wide
collapse, where per-candidate ranking measures fragility ordering within
an avalanche. Secondary metrics (AUROC, Overlap@$K_p$, tail $\rho$) for
the two headline probes appear in the appendix.

\begin{table*}[t]
\caption{\textbf{Channel-conditioned prediction of realized damage
(TOFU tofu-a180, sealed audit fold, $n{=}300$, Qwen2.5-1.5B fp32).}
Spearman $\rho$ between each sealed pre-unlearning score and realized
per-candidate damage, per objective; objectives grouped by declared
damage channel, predictors by family; column-best in bold. ``Bwd-free''
marks predictors needing no candidate-side backward pass. The
preregistered interaction $\Delta$ (Eq.~\eqref{eq:interaction}) between
the headline probes is positive with bootstrap CIs excluding zero on
both declared pairs: GradDiff/RMU $\Delta=+0.300$, 95\% CI
$[+0.131,+0.466]$; NPO/RepNoise $\Delta=+0.470$, 95\% CI
$[+0.232,+0.698]$ (2{,}000 candidate bootstraps). Cells are
single-request point estimates; seed replication is reported in the
appendix. $^{\dag}$did not reach the preregistered forget criterion
within budget (weak-signal column, interpret with care);
$^{\ddag}$reached the criterion but with model-wide utility collapse.}
\label{tab:channel-matrix}
\centering
\scriptsize
\setlength{\tabcolsep}{2.6pt}
\begin{tabular}{@{}llc cccccc ccc cc@{}}
\toprule
& & & \multicolumn{6}{c}{\textbf{Output / loss-gradient channel}}
& \multicolumn{3}{c}{\textbf{Representation channel}}
& \multicolumn{2}{c}{Channel mean} \\
\cmidrule(lr){4-9}\cmidrule(lr){10-12}\cmidrule(l){13-14}
Family & Predictor & Bwd-free & GradDiff & NPO & SimNPO & GRU &
GA$^{\ddag}$ & IdkDPO$^{\dag}$ & RMU & RepNoise & CB$^{\dag}$ & Out & Rep \\
\midrule
Gradient magnitude & Exact gradient norm & \xmark &
0.40 & 0.39 & 0.39 & 0.43 & 0.20 & 0.41 & -0.04 & -0.16 & -0.23 & 0.37 & -0.14 \\
& Randomized sensitivity (\textbf{ours}) & \cmark &
0.41 & \textbf{0.40} & \textbf{0.40} & \textbf{0.43} & 0.20 & \textbf{0.41} & -0.05 & -0.17 & -0.25 & 0.38 & -0.15 \\
\addlinespace
Representation proximity & Hidden-state kNN & \cmark &
\textbf{0.49} & 0.23 & 0.21 & 0.42 & \textbf{0.50} & 0.05 & 0.34 & 0.14 & \textbf{0.09} & 0.32 & 0.19 \\
& Sentence-embedding kNN & \cmark &
0.28 & 0.15 & 0.14 & 0.25 & 0.35 & 0.10 & 0.10 & \textbf{0.21} & 0.05 & 0.21 & 0.12 \\
& Lexical overlap & \cmark &
0.39 & 0.17 & 0.16 & 0.34 & 0.18 & 0.00 & \textbf{0.39} & -0.02 & 0.05 & 0.21 & 0.14 \\
\addlinespace
Alignment (rejected) & Forget-direction FD & \cmark &
-0.06 & -0.08 & -0.07 & -0.06 & 0.07 & -0.09 & -0.11 & 0.11 & -0.06 & -0.05 & -0.02 \\
\addlinespace
Controls & Random ranking & \cmark &
0.02 & 0.10 & 0.10 & 0.03 & 0.01 & 0.13 & -0.05 & -0.11 & -0.00 & 0.07 & -0.06 \\
\bottomrule
\end{tabular}
\end{table*}


\subsection{Main Result II: Crossed Channel-Matched Protection}
\label{sec:exp-protection}

Prediction is operationally useful only if routing a budget through it
improves the sealed audit. Table~\ref{tab:crossed} reports the crossed
design: each parent (GradDiff for the output channel, RMU for the
representation channel) runs unmodified to its criterion-reaching
checkpoint, then each selector arm spends an identical repair budget under
the identical guarded procedure of
Section~\ref{sec:method-protection}; only the ranking that fills the
protect pool differs. The preregistered success criterion is that the
matched arm attains the lowest audit CVaR for \emph{both} parents. The
exact-norm selector bounds what the gradient family could buy at any cost;
comparing it with randomized sensitivity prices the backward-free
discount. Because the guard caps forgetting regression, matched protection
is claimed at matched forgetting -- the reach and forget-recall columns
verify that no arm buys retention by under-forgetting.

\begin{table*}[t]
\caption{\textbf{Crossed protection: parent objective $\times$ protection
selector (TOFU, sealed native audit).} Every arm of a parent starts from
the same reaching checkpoint and uses the same repair budget, projection,
guard, and audit; only the selector changes. Match labels follow the
preregistered channel table. Success = matched CVaR lowest per parent.
Conditional cells report $n_{\mathrm{reach}}/n_{\mathrm{total}}$ and
estimate [95\% CI]; lower damage is better.}
\label{tab:crossed}
\centering
\scriptsize
\setlength{\tabcolsep}{3.0pt}
\begin{tabular}{@{}lll cccccc@{}}
\toprule
Parent (channel) & Selector & Match & Reach $\uparrow$ &
Forget recall $\downarrow$ & Para.\ recall $\downarrow$ &
Mean $\Delta$NLL $\downarrow$ & CVaR$_{.95}$ $\downarrow$ &
Utility ret.\ $\uparrow$ \\
\midrule
GradDiff (output) & none (parent alone) & --- &
\tblph & \tblph & \tblph & \tblph & \tblph & \tblph \\
& random & random &
\tblph & \tblph & \tblph & \tblph & \tblph & \tblph \\
& hidden-state proximity & mismatched &
\tblph & \tblph & \tblph & \tblph & \tblph & \tblph \\
& randomized sensitivity (ours) & \textbf{matched} &
\tblph & \tblph & \tblph & \tblph & \tblph & \tblph \\
& exact gradient norm & matched (ceiling) &
\tblph & \tblph & \tblph & \tblph & \tblph & \tblph \\
\addlinespace
RMU (representation) & none (parent alone) & --- &
\tblph & \tblph & \tblph & \tblph & \tblph & \tblph \\
& random & random &
\tblph & \tblph & \tblph & \tblph & \tblph & \tblph \\
& randomized sensitivity & mismatched &
\tblph & \tblph & \tblph & \tblph & \tblph & \tblph \\
& hidden-state proximity & \textbf{matched} &
\tblph & \tblph & \tblph & \tblph & \tblph & \tblph \\
\bottomrule
\end{tabular}
\end{table*}


\subsection{Backward-Free Fidelity and Cost}
\label{sec:exp-fidelity}

Table~\ref{tab:bwfree} substantiates the probe's two supporting claims.
Fidelity: within the gate-certified regime, randomized sensitivity ranks
candidates in near-perfect agreement with the exact gradient norm, and its
damage prediction retains most of the exact profile's power; outside the
regime (radius in the cancellation zone, or low-precision underflow) the
gate correctly withholds certification -- we report both sides of the
boundary. Cost: profiling an entire candidate universe takes $2R$ batched
forward sweeps at inference memory, versus per-candidate backward passes
at training memory; wall-clock and peak-memory measurements at 7B compare
the randomized probe against exact-norm, JVP, and vectorized
gradient-dot implementations. The advantage we claim is memory and
backward-freedom, not universal wall-clock supremacy.

\begin{table*}[t]
\caption{\textbf{The backward-free probe: fidelity to the exact profile
and profiling cost.} Left: rank agreement with the exact gradient norm
inside and outside the preregistered fidelity gate (Qwen models, fp32);
at the gate-certified setting the probe's damage prediction matches the
exact profile on all nine objectives of Table~\ref{tab:channel-matrix}
to within $\pm 0.01$. Right: measured wall-clock and peak memory per
symmetric direction pair over one candidate universe (Qwen2.5-7B, bf16,
batch 4); randomized sensitivity runs $R$ such pairs at the same
$18.3$\,GB peak -- time scales linearly in $R$ while memory stays at
inference level, so the claim is memory and backward-freedom, not
universal wall-clock supremacy.}
\label{tab:bwfree}
\centering
\scriptsize
\setlength{\tabcolsep}{3.0pt}
\begin{tabular}{@{}lcc c lcc@{}}
\toprule
\multicolumn{3}{c}{\textbf{Fidelity (fp32)}} & &
\multicolumn{3}{c}{\textbf{Cost (7B, bf16, per pass)}} \\
\cmidrule(r){1-3}\cmidrule(l){5-7}
Regime & $\rho$(rand, exact) & Damage-$\rho$ retained & &
Profiler & Time (s) & Peak mem (GB) \\
\midrule
Gate-certified ($\eta{=}3\mathrm{e}{-}3$, $R{=}64$) & 0.94 &
9/9 obj.\ within $\pm 0.01$ & &
Symmetric FD pair ($2\times$ fwd sweep) & 5.97 & \textbf{18.3} \\
Cancellation zone ($\eta{=}3\mathrm{e}{-}4$) & 0.50 & \tblph & &
Exact JVP (forward-mode) & 9.9 & 27.4 \\
Low-precision underflow (bf16) & \tblph & \tblph & &
Vectorized gradient-dot & 10.0 & 35.2 \\
Curvature zone (large $\eta$) & \tblph & \tblph & &
Streaming candidate backward & 17.4 & 18.7 \\
\bottomrule
\end{tabular}
\end{table*}


\subsection{Soft Routing: The Mixture Sweep}
\label{sec:exp-mixture}

The mixture score of Eq.~\eqref{eq:mixture} turns the discrete router
into a measurable one-parameter family, and its two sides have very
different costs. The \emph{prediction} sweep is free: both headline
scores are already sealed per candidate, so
$\rho(s_\alpha,\text{damage})$ over a grid
$\alpha\in\{0,0.125,\ldots,1\}$ is pure post-processing of existing
artifacts for every objective at once. The channel prediction is a pair
of mirrored tendencies -- output-channel objectives improving toward
high $\alpha$, representation-channel objectives toward low $\alpha$ --
with interior optima exactly where the spectrum clause predicts residual
off-channel signal (e.g.\ GradDiff on this proximity-rich benchmark,
where the proximity probe retains substantial $\rho$).
\resph{Per-objective $\alpha$ curves and prediction-optimal
$\alpha^{*}_{\mathrm{pred}}$ per channel.}

The \emph{protection} sweep reruns the crossed pipeline with the protect
pool $\operatorname{TopK}$ by $s_\alpha$ for
$\alpha\in\{0,0.25,0.5,0.75,1\}$, holding parent, budget, and repair
fixed; the endpoints coincide with the matched/mismatched arms of
Table~\ref{tab:crossed}, so the sweep interpolates between rows that are
already measured. The quantity of interest is the gap between
$\alpha^{*}_{\mathrm{pred}}$ and $\alpha^{*}_{\mathrm{prot}}$ per
parent: a systematic shift of $\alpha^{*}_{\mathrm{prot}}$ toward the
proximity end would confirm that repair generalization propagates
through representation-space neighborhoods independently of the parent's
damage channel -- the mechanism suggested by the crossed observations --
and would refine the routing recommendation from ``match the channel''
to ``match the channel for prediction; weight proximity for repair
transfer.'' \resph{Protection $\alpha$ curves per parent;
$\alpha^{*}_{\mathrm{pred}}$ vs $\alpha^{*}_{\mathrm{prot}}$.}

\subsection{Replication, Stress, and Positioning}
\label{sec:exp-supporting}

Replication runs repeat the reduced crossed roster on WMDP-bio, MUSE-News,
and RWKU; the confirmatory quantity in each is the interaction, not any
absolute correlation, since each dataset shifts the probes' base rates.
Stress runs target the honesty clauses: PISTOL maximizes forget--retain
entanglement (proximity probes risk over-selection), and MUSE-Books tests
verbatim-memorization removal. Seed and request replication on TOFU, the
specificity controls (benign SFT, unrelated requests, norm-matched random
motion), repair-budget sweeps, and the strong-method comparison that
places the full pipeline alongside recent unlearners under the same
contract are reported in the appendix. The pipeline is a layer over
parents, so that comparison is context, not the claim: the claim is the
pair of interactions, which no channel-blind method -- however strong --
can exhibit.

\section{Discussion}
\label{sec:implications}

\paragraph{From a single score to a conditional one.}
The literature's implicit model has been
$\text{risk}(x)=f(\text{model},\text{data})$: some one signal --
similarity, alignment, magnitude -- should rank vulnerable candidates for
any unlearner. Our results replace it with
$\text{risk}(x \mid \text{channel})$. The interaction is the right
evidence for this claim precisely because it is immune to the two
confounds that plague probe leaderboards: a uniformly better probe and a
uniformly more predictable objective both cancel in the double
difference. The asymmetry we observe -- proximity probes retain partial
signal on output-channel objectives while gradient probes go essentially
blind on RMU -- also explains why single-score studies reached
inconsistent conclusions: their conclusions depended on which objectives
happened to be in the roster. That the conditioning variable is trivial
to evaluate -- a static read of the loss definition -- is a feature, not
a limitation: the scientific content lies not in the taxonomy but in the
demonstrated law that this five-second label governs which predictor
family carries signal, by how much, and with which failure mode
(gradient probes do not merely weaken on the representation channel;
they go blind). A simple declared predicate with a non-obvious,
falsifiable consequence is precisely what makes the router deployable
and the claim preregistrable.

\paragraph{Operational picture.}
A practitioner declares the checkpoint, candidate scope, parameter block,
loss, and probe settings; reads the parent objective's channel off its
loss definition; and runs the matched probe. The profile's high-risk tail
localizes where post-unlearning measurement should concentrate, and it
allocates the protection budget before any audit is opened. The router
adds no tuning surface: the channel table is a lookup, and both probes are
forward-only at candidate scale. The profile does not decide whether a
request should be served by unlearning, retraining, or refusal, and a high
score does not mandate preservation; those remain application decisions.

\paragraph{Selection is not confirmation.}
A profile-selected pool may guide optimization, but improvement on that
pool is not evidence of preservation -- it rewards the method for
protecting what it was handed. All confirmatory outcomes are therefore
sealed-audit quantities, with both mean and tail reported because
concentrated damage hides in means.

\paragraph{What the score means.}
Each profile is conditional on the model, block, loss, precision, and
radius, and should be recomputed after material model changes. Randomized
sensitivity is an update-sensitivity prior rather than an
unlearning-specific quantity; its unlearning reading is licensed by the
specificity controls, and its trustworthiness by the fidelity gate rather
than by assumption.

\paragraph{Prediction and intervention are separate tests.}
A predictor can rank damage yet fail to support safe repair, and a repair
procedure can help for reasons unrelated to its selector. The crossed
design keeps the two claims apart: Table~\ref{tab:channel-matrix}
establishes conditional prediction, and Table~\ref{tab:crossed}
establishes that acting on the matched prediction -- and only the matched
prediction -- moves the sealed tail.


\section{Limitations}
\label{sec:limitations}

The channel dichotomy is an empirical classification of where current
objectives read the model, not a law: hybrid objectives inherit partial
signal from both families, and an objective reading a third quantity would
require a third probe. We preregistered this clause and stress-test it
(PISTOL, MUSE-Books) rather than treating the two channels as exhaustive.

Both probes are local and block-scoped: profiles are conditional on the
declared block, loss, precision, and radius, certified by the fidelity
gate only within that geometry, and should be recomputed after material
model changes. The randomized probe is additionally an update-sensitivity
prior whose unlearning-specific reading rests on the reported specificity
controls.

The profile ranks an enumerated candidate universe and a finite sealed
audit; behavior outside both is uncovered, and shared templates across
folds can shrink the effective sample. Guarded repair caps forgetting
regression empirically -- first-order projection plus a paired guard -- and
certifies neither data removal nor robustness to extraction or relearning
attacks. Claims beyond the completed datasets, scales, and request types
remain bounded by the transfer experiments.


\section{Conclusion}
\label{sec:conclusion}

This paper treats retain-side susceptibility as conditional on where an
unlearning objective places its forget-suppression signal. Its mathematical
contribution is one backward-free identity: squared loss responses to
shared symmetric random parameter shakes recover candidate gradient energy
up to finite-difference error. Output-readout objectives use that exposure
profile; representation-readout objectives use hidden-state proximity; the
confirmatory router remains the fixed $0/1$ endpoint rule.

The empirical claim is deliberately falsifiable: the frozen family-level
interaction must be positive, and matched allocation must lower sealed-audit
tail damage for each crossed parent at comparable forgetting. Interior
$s_\alpha$ values only diagnose cross-channel signal and never rescue a
failed endpoint test. \resph{Insert the final interaction, fidelity, and
crossed-protection outcomes after the sealed runs.} The resulting layer is
a profiler and protection audit around existing unlearners, not a new
unlearning objective or a universal susceptibility score.


%% ==========================================================================
%% Acknowledgments are omitted in the anonymous submission.

\bibliographystyle{ACM-Reference-Format}
\bibliography{references}

%% ==========================================================================
\appendix
% The appendix contains reproducibility detail and secondary analyses only.
% Claim decisions and their decisive bounds remain in the main paper.

\section{Reproducibility and Statistical Details}
\label{app:evaluation-contract}
\label{app:datasets}
\label{app:objectives}

The released campaign manifest is the source of truth for every request,
candidate, semantic group, checkpoint, seed, random pool, and configuration
hash. Table~\ref{tab:datasets} is a readable digest, not a replacement for
that manifest. The denominator is the predeclared matrix: planned rows remain
visible when they are never attempted, terminate early, fail integrity, do not
reach the common forgetting criterion, or lack complete all-arm support.

\begin{table*}[t]
\caption{\textbf{Compact campaign and outcome contract.} Panel A records the
predeclared scope. The three development folds are mutually disjoint and
target-free. Panel B defines constraint-matched forgetting and the common
utility floor; exact metrics and thresholds are frozen in the manifest before
target outcomes.}
\label{tab:datasets}
\label{tab:objective-config}
\label{tab:constraint-contract}
\centering
\scriptsize
\setlength{\tabcolsep}{2.7pt}
\resizebox{\textwidth}{!}{%
\begin{tabular}{@{}llllcccc@{}}
\toprule
\multicolumn{8}{@{}l}{\textbf{A. Predeclared setting scope}} \\
Dataset / role & Model / precision & $\mathcal D_{\rm cal}/\mathcal D_{\rm pred}/\mathcal D_{\rm prot}$ &
Target requests & Primary parents & Seeds & $|\mathcal C_{\rm disc}|/|\mathcal C_{\rm audit}|$ & Planned cells \\
\midrule
TOFU / primary & Qwen2.5-1.5B / fp32 & \tblph & \tblph & 7 & \tblph & \tblph & \tblph \\
TOFU / scale & Qwen2.5-7B / fp32 & \tblph & \tblph & 7 & \tblph & \tblph & \tblph \\
TOFU / family & Llama-3.1-8B / fp32 & \tblph & \tblph & 7 & \tblph & \tblph & \tblph \\
WMDP-bio / MMLU & \tblph & \tblph & \tblph & 7 & \tblph & \tblph & \tblph \\
MUSE-News & \tblph & \tblph & \tblph & 7 & \tblph & \tblph & \tblph \\
RWKU & \tblph & \tblph & \tblph & 7 & \tblph & \tblph & \tblph \\
MUSE-Books / stress & \tblph & \tblph & \tblph & 7 & \tblph & \tblph & \tblph \\
PISTOL / stress & \tblph & \tblph & \tblph & 7 & \tblph & \tblph & \tblph \\
\bottomrule
\end{tabular}}

\vspace{3pt}
\resizebox{0.88\textwidth}{!}{%
\begin{tabular}{@{}llllll@{}}
\toprule
\multicolumn{6}{@{}l}{\textbf{B. Shared final-checkpoint feasibility region}} \\
Audit & Metric & Favorable direction & Frozen boundary & Stopping role & Reported slack \\
\midrule
Direct forgetting & \tblph & \tblph & \tblph & parent entry + final constraint & yes \\
Paraphrase forgetting & \tblph & \tblph & \tblph & final constraint & yes \\
Extraction / generation & greedy autoregressive gold-token recall & lower & \tblph & final constraint & yes \\
General utility & \tblph & \tblph & \tblph & final floor & yes \\
Native retain audit & mean/CVaR$_{.95}$ $\Delta$NLL + dataset-native metric & lower/higher & frozen NI margin & outcome only & n.a. \\
\bottomrule
\end{tabular}}
\end{table*}

The primary objective roster is GradDiff, NPO, SimNPO, GRU, RMU, RepNoise,
and Circuit Breakers; GA and IdkDPO are stress controls. For each parent the
artifact records the implementation identifier, source commit, trainable
block, learning rate, update and processed-token budgets, saved-step grid,
forget/reference coefficients, stopping rule, and hashes of all frozen
inputs. Readout labels are descriptive and never select a mixture endpoint.

The candidate manifest freezes $\mathcal C$, the discovery/audit split by
semantic group, repair eligibility, and a score-independent random stream. Direct
answer duplicates, forget-fact paraphrases, and template-derived restatements
are ineligible for repair; entity overlap is retained as a disclosed stratum.
Only $\mathcal C_{\rm disc}$ fits empirical midrank maps or supplies repair
examples. Audit scores are computed at $\theta_0$ and sealed. Candidate IDs
break ties only when selecting the last member of a top-$K_p$ pool.
The neutral stream and differentiable guard sets are selector-independent.
General-utility checkpoint selection uses a separate frozen set; neither
$\mathcal C_{\rm audit}$ nor its native retained outcome is used by a guard,
stopping rule, or checkpoint selector.

\subsection{Units, funnels, and decisions}
\label{app:outcomes}

Every prediction correlation is formed within one
model--request--seed--parent cell. Seeds are averaged within requests and
requests receive equal weight. The paired hierarchical bootstrap resamples
requests, trajectory seeds within requests, and semantic groups within
requests. Its replicate count, seed, and
finite-sample percentile rule are manifest fields. No candidate is pooled
across requests before a correlation is computed.

Within a request, empirical $\mathrm{CVaR}_{.95}$ is the mean of the largest
5\% damage values, using fractional weight at the quantile boundary. When an
audit has fewer than 20 candidates, this reduces to the maximum and the row is
flagged as low-tail-support rather than pooled with larger audits.

The audit exposes four denominators. $N_{\rm prof}$ counts all predeclared
profiles and $N_{\rm valid}$ those passing the outcome-blind integrity check.
$N_{\rm att}$ counts all predeclared parent trajectories and $N_{\rm reach}$
those reaching the common entry criterion. $N_{\rm rv}$ is the
reached--valid intersection. Prediction common support requires joint score,
both endpoints, the frozen simple control, and damage on the same units. A
tail cell is eligible only when its positive-damage count is at least $m=K_p$,
with $q=m/|\mathcal C_{\rm audit}|$. Tail-ineligible cells remain in
the RQ1 correlation denominator; $N_{\rm tail}/N_{\rm rv}$ is always reported,
and the tail claim requires at least 0.80 coverage. Protection feasibility
requires mixture, no repair, every frozen random draw, $S_0$, and $S_1$ to
satisfy all four constraints on the same units; common support additionally
requires both damage outcomes and the native retained metric for every arm.
No missing outcome is imputed.

RQ1 is the intersection--union test (IUT)
$\rho_S>0\land g_G>0\land g_H>0\land L_{\rm tail}>0$, with the tail term
additionally requiring at least 0.80 coverage. RQ2 requires
$f_\rho>0.80$, $f_K>0.70$, $g_H>0$, and $g_{\rm ctl}>0$ through their
one-sided lower bounds. Protection is the IUT over four comparators, two
$\Delta$NLL outcomes, and one native-metric non-inferiority margin per
comparator. Thus success requires four RQ1 lower bounds, four RQ2 lower
bounds, and eight superiority plus four non-inferiority bounds for RQ3 to
cross their frozen boundaries in the favorable direction. Statistical
success and eligibility are separate
fields. A fallback weight, invalid profile, missing
expected cell, non-reach, infeasible arm, incomplete random-draw roster, or
incomplete common support forces claim pass to ``no,'' even if a conditional
effect is printable. Exact-gradient allocation is a same-estimand numerical
reference outside the protection IUT.

Parent rows use their one-sided IUT level without a within-row multiplicity
penalty. The later setting-level rule can select a passing parent from each
readout group, so that selection applies a Bonferroni threshold across the four
output-readout or three representation-readout parents. This correction is
used only for the breadth statement; it does not hide the complete row roster.

The paper-level breadth rule is defined in Section~\ref{sec:exp-supporting}; it
does not pool parent rows or turn the least-favorable row in
Table~\ref{tab:robustness} into a simultaneous confidence bound.

\FloatBarrier

\section{Additional Prediction Analyses}
\label{app:prediction-supp}

The main paper defines the tail premise and carries every parent-level claim
decision. Table~\ref{tab:tail-structure} provides the parent breakdown and
secondary controls. Tail support is conditioned on a realized parent damage
vector, not on profile validity, so a failed profiler cannot erase evidence
about whether aggregate retention hides a coherent tail. The permutation
keeps request-specific group sizes fixed. If a request has zero positive
damage mass, its concentration ratio is undefined: we report it as a
zero-mass request and exclude it only from that ratio, never from the
prediction denominator. Tail size and minimum support are frozen in the
manifest.

\begin{table*}[t]
\caption{\textbf{Additional tail, reference, and construction analyses.}
Panel A expands the premise behind Figure~\ref{fig:channel-main}A. Panel B
compares pre-unlearning references on equal parent/request support. Panel C
contains representative one-factor checks; the full grid and bootstrap draws
are artifacts. None of these panels replaces the main table's absolute and
two-endpoint prediction test.}
\label{tab:tail-structure}
\label{tab:prospective-prediction}
\label{tab:profile-construction}
\label{tab:horizon-sensitivity}
\centering
\scriptsize
\setlength{\tabcolsep}{2pt}
\resizebox{0.88\textwidth}{!}{%
\begin{tabular}{@{}lccccccc@{}}
\toprule
\multicolumn{8}{@{}l}{\textbf{A. Damage-tail evidence on repair-held-out audit behaviors}} \\
Parent & Mean damage & CVaR$_{.95}$ & $M_q/q$ [95\% CI] &
$L_{\rm grp}$ [95\% CI] & Permutation $p$ & Damage support $n/N$ & Tail $n$ \\
\midrule
GradDiff & \tblph & \tblph & \tblph & \tblph & \tblph & \tblph & \tblph \\
NPO & \tblph & \tblph & \tblph & \tblph & \tblph & \tblph & \tblph \\
SimNPO & \tblph & \tblph & \tblph & \tblph & \tblph & \tblph & \tblph \\
GRU & \tblph & \tblph & \tblph & \tblph & \tblph & \tblph & \tblph \\
RMU & \tblph & \tblph & \tblph & \tblph & \tblph & \tblph & \tblph \\
RepNoise & \tblph & \tblph & \tblph & \tblph & \tblph & \tblph & \tblph \\
CB & \tblph & \tblph & \tblph & \tblph & \tblph & \tblph & \tblph \\
\bottomrule
\end{tabular}}

\vspace{3pt}
\resizebox{\textwidth}{!}{%
\begin{tabular}{@{}llccccccc@{}}
\toprule
\multicolumn{9}{@{}l}{\textbf{B. Pre-unlearning reference roster; descriptive equal-parent summaries}} \\
Family & Predictor & Candidate reverse & All $\rho$ & Output $\rho$ &
Representation $\rho$ & Top-$q$ recall & Common support & Role \\
\midrule
Gradient & exact gradient energy & yes & \tblph & \tblph & \tblph & \tblph & \tblph & numerical reference \\
Directional & forget-direction symmetric FD & no & \tblph & \tblph & \tblph & \tblph & \tblph & baseline \\
Loss response & one random direction & no & \tblph & \tblph & \tblph & \tblph & \tblph & baseline \\
Representation & sentence-embedding proximity & no & \tblph & \tblph & \tblph & \tblph & \tblph & baseline \\
Lexical & token-overlap proximity & no & \tblph & \tblph & \tblph & \tblph & \tblph & baseline \\
Representation & request-shuffled proximity & no & \tblph & \tblph & \tblph & \tblph & \tblph & negative control \\
Loss & initial NLL & no & \tblph & \tblph & \tblph & \tblph & \tblph & baseline \\
Length & answer length & no & \tblph & \tblph & \tblph & \tblph & \tblph & baseline \\
Random & repeated frozen rankings & no & \tblph & \tblph & \tblph & \tblph & \tblph & chance \\
\bottomrule
\end{tabular}}

\vspace{3pt}
\resizebox{\textwidth}{!}{%
\begin{tabular}{@{}lllccccc@{}}
\toprule
\multicolumn{8}{@{}l}{\textbf{C. Representative one-factor checks}} \\
Axis & Primary / variant & Scope & Joint $\rho$ & min gain [min LB] &
Top-$q$ recall & Common support & Eligible/pass \\
\midrule
Representation layer & final / penultimate & one output + one representation parent & \tblph & \tblph & \tblph & \tblph & \tblph \\
Pooling & answer-token mean / final token & same & \tblph & \tblph & \tblph & \tblph & \tblph \\
Loss-shake block & frozen late MLP / last layer & same & \tblph & \tblph & \tblph & \tblph & \tblph \\
Mixture & parent-specific / one global weight & all primary parents & \tblph & \tblph & \tblph & \tblph & \tblph \\
Horizon & first reach / fixed or terminal & all reached parents & \tblph & \tblph & \tblph & \tblph & \tblph \\
\bottomrule
\end{tabular}}
\end{table*}

All reference predictors use the same candidate folds and outcome horizons.
Construction ablations keep the primary prediction weight fixed; a separately
reported development refit diagnoses weight mismatch but cannot replace that
ablation. The request-shuffled control draws a different forget request of the
same declared size before target outcomes.

\FloatBarrier

\section{Loss-Shake Numerical Validation}
\label{app:estimator-stability}
\label{app:probe-details}
\label{app:randomized-impl}

For each Gaussian draw, the implementation normalizes
$v_r=g_r/\lVert g_r\rVert_2$, saves the declared parameter block, applies
$\theta_0+\eta v_r$ and $\theta_0-\eta v_r$ from that same copy, and restores
it before the next direction. Candidate order,
batching, masks, padding, precision, and dropout state are identical within a
pair. Signed responses are accumulated in fp32 and persisted before squaring,
so nested-$R$, direction-bank, and split-half analyses require no new model
evaluation. The dimension factor in Eq.~\eqref{eq:loss-shake-identity} is
stored even though it does not affect ranking.

The confirmatory operating point uses fp32 parameters. Bfloat16 is retained
only as a numerical boundary check: a configuration is invalid when the
realized block displacement or changed-coordinate fraction falls below its
frozen perturbation-survival threshold.

The exact reference streams one candidate at a time, records its gradient
energy on the same block, and discards the vector. Outcome-blind development
calibration chooses the smallest $(R,\eta)$ satisfying the frozen criteria for
exact-rank fidelity, top-$K_p$ overlap, perturbation survival, and
direction-split stability. A
target profile is valid only when the latter forward-only checks pass on
$\mathcal C_{\rm disc}$ before the sealed audit is opened.
\label{app:fidelity-gate}

\begin{table*}[t]
\caption{\textbf{Loss-shake fidelity--cost frontier and numerical
sensitivity.} Timing uses device synchronization and a common measurement
boundary; memory is peak allocated device memory. The selected operating
point is marked frozen. Exact gradients are evaluated only as a reference.}
\label{tab:bwfree}
\label{tab:estimator-validation}
\centering
\scriptsize
\setlength{\tabcolsep}{2.2pt}
\resizebox{\textwidth}{!}{%
\begin{tabular}{@{}lllcccccc@{}}
\toprule
\multicolumn{9}{@{}l}{\textbf{A. Matched fidelity--cost frontier}} \\
Model / block / precision & Profiler & $R$ & $\rho$ vs. exact & overlap@$K_p$ &
Pool passes / candidate reverse & Time & Peak memory & Integrity valid $n/N$ \\
\midrule
Qwen2.5-1.5B / \tblph / fp32 & exact energy & --- & 1.00 & 1.00 & candidate F+B / yes & \tblph & \tblph & n.a. \\
& exact energy (batched, if feasible) & --- & 1.00 & 1.00 & batched F+B / yes & \tblph & \tblph & n.a. \\
& loss-shake & \tblph & \tblph & \tblph & $2R$ pool F / no & \tblph & \tblph & \tblph \\
& loss-shake (frozen) & \tblph & \tblph & \tblph & $2R$ pool F / no & \tblph & \tblph & \tblph \\
Qwen2.5-7B / \tblph / fp32 & exact energy & --- & 1.00 & 1.00 & candidate F+B / yes & \tblph & \tblph & n.a. \\
& loss-shake (frozen) & \tblph & \tblph & \tblph & $2R$ pool F / no & \tblph & \tblph & \tblph \\
Llama-3.1-8B / \tblph / fp32 & loss-shake (frozen) & \tblph & \tblph & \tblph & $2R$ pool F / no & \tblph & \tblph & \tblph \\
\bottomrule
\end{tabular}}

\vspace{3pt}
\resizebox{0.9\textwidth}{!}{%
\begin{tabular}{@{}llccccc@{}}
\toprule
\multicolumn{7}{@{}l}{\textbf{B. One-factor numerical sensitivity}} \\
Axis & Frozen / variants & $\rho$ range & overlap range & split-half range &
Valid $n/N$ & Boundary diagnosis \\
\midrule
Directions & $R$ / nested prefixes and new banks & \tblph & \tblph & \tblph & \tblph & \tblph \\
Radius & $\eta$ / log-grid neighbors & \tblph & \tblph & \tblph & \tblph & \tblph \\
Precision & declared / fp32, bf16 & \tblph & \tblph & \tblph & \tblph & \tblph \\
Pool size & full / nested pools, fixed common subset & \tblph & \tblph & \tblph & \tblph & \tblph \\
\bottomrule
\end{tabular}}
\end{table*}

Rank agreement validates computation of $\widehat q_G$ only. Fidelity alone
establishes neither component's predictive added value; the main paired gains
and simple-control comparison decide that separate claim.

\FloatBarrier

\section{Protection Sensitivity and Boundary Checks}
\label{app:protection-supp}
\label{app:protection-details}

\subsection{Common repair harness}
\label{app:guard-proof}
\label{app:repair-protocol}

Every active arm begins from the same first criterion-reaching parent
checkpoint. It selects exactly $K_p$ eligible discovery examples, but shares
the neutral stream $\mathcal R_0$, example order, optimizer random stream,
processed-token budget, projection, guards, and save schedule. The neutral
stream enters $L_{\rm rep}$ through the entry-anchored KL term in
Eq.~\eqref{eq:budgeted-repair}; no susceptibility score selects it. The token
budget charges protected, neutral, constraint-refresh, guard, and
checkpoint-evaluation tokens.

Projection covers only differentiable surrogate constraints. Each is oriented
as $c_j(\theta)\leq0$ on a frozen guard set and is active when
$c_j(\theta)\geq-\kappa_j$. The matrix $G$ is refreshed every $m_{\rm ref}$
accepted steps and Eq.~\eqref{eq:repair-projection} uses ridge $\lambda$.
For a forget token $(i,s)$, the anchored adverse margin is
\[
 r^{f}_{is}(\theta)=\log p_\theta(y_{is}\mid y_{i,<s})-
 \log p_{t^\dagger}(y_{is}\mid y_{i,<s});
\]
for a neutral or utility token it is the NLL increase
\[
 r^{u}_{is}(\theta)=-\log p_\theta(y_{is}\mid y_{i,<s})+
 \log p_{t^\dagger}(y_{is}\mid y_{i,<s}).
\]
A proposed step must keep every token margin below its frozen
$\epsilon^{\rm tok}$ and every example-average margin below
$\epsilon^{\rm ex}$. Otherwise the parameters are restored, the step size is
halved, and the step is retried up to $J_{\rm retry}$ times; exhaustion stops
the arm. The manifest freezes $\beta$, $\kappa_j$, $m_{\rm ref}$, $\lambda$,
both tolerances, $J_{\rm retry}$, and guard IDs before target runs.

Non-differentiable extraction, generation, and full utility metrics are not
represented in $G$. They are evaluated only at the shared save schedule on
selector-safe sets. The reported output is the last checkpoint satisfying
every row of Table~\ref{tab:constraint-contract}; if none does, the arm is
infeasible. Sealed $\Delta$NLL and native retained outcomes are opened only
after this selection. The artifact records refreshes, retries, accepted and
rejected updates, margins, and the selected step.

Repeated random allocation uses $B_{\rm rand}$ pools frozen before outcomes.
Every draw must complete and satisfy the same constraints. Draw IDs are
resampled and averaged inside their trajectory within each bootstrap
replicate. Diagnostic arms replace only the selector: exact-energy allocation
substitutes $q_G^\star$ for $\widehat q_G$; prediction-weight allocation uses
$\widehat\alpha_{\rm pred}$; uniform rehearsal ignores ranking; and a
realized-discovery-damage oracle gives an unattainable repair-capacity
reference. All remain outside the four-comparator claim.

\subsection{Budget and specificity checks}
\label{app:specificity}

The specificity analysis keeps the sealed target profile fixed. Other-request
motion uses a size-matched forget set; random motion matches the target block
displacement and uses directions independent of profiling. These controls
bound interpretation: loss-shake energy is intentionally direction- and
request-agnostic, whereas proximity should carry the request-specific part.
MUSE-Books and PISTOL remain boundary settings and never count toward the
paper-level breadth rule.

Hyperparameter grids, source-to-target weight carryover, exploratory traces of
directional persistence, and early hypothesis development are artifact
diagnostics rather than paper tables. They never select a target rule or fill
a confirmatory cell.

\begin{table*}[!t]
\caption{\textbf{Protection sensitivity and specificity boundaries.} Panel A
varies the protected fraction under the unchanged five-arm contract. Panel B
asks whether each frozen component responds specifically to the target
request. Primary parent decisions and their decisive worst UCB remain in
Table~\ref{tab:core-evidence}; complete arm effects, constraint vectors, guard
events, and random-pool rows are released with the artifact.}
\label{tab:crossed}
\label{tab:budget-sweep}
\label{tab:protection-ablations}
\label{app:budget-sweep}
\centering
\scriptsize
\setlength{\tabcolsep}{1.8pt}
\resizebox{0.72\textwidth}{!}{%
\begin{tabular}{@{}lccccc@{}}
\toprule
\multicolumn{6}{@{}l}{\textbf{A. Fixed-budget sensitivity}} \\
Budget & Eligible/pass & worst $\Delta$ [UCB] & Bottleneck &
min F/U slack & Updates/support \\
\midrule
0.05 & \tblph & \tblph & \tblph & \tblph/\tblph & \tblph/\tblph \\
0.10 (primary) & \tblph & \tblph & \tblph & \tblph/\tblph & \tblph/\tblph \\
0.20 & \tblph & \tblph & \tblph & \tblph/\tblph & \tblph/\tblph \\
\bottomrule
\end{tabular}}

\vspace{3pt}
\resizebox{0.72\textwidth}{!}{%
\begin{tabular}{@{}llcccc@{}}
\toprule
\multicolumn{6}{@{}l}{\textbf{B. Request specificity and motion controls}} \\
Motion & Conditioning & $\Delta\rho_{G/H/S}$ & top-$q$ lift & Disp. matched & Support \\
\midrule
Target unlearning & target request & \tblph/\tblph/\tblph & \tblph & yes & \tblph \\
Other-request & different request & \tblph/\tblph/\tblph & \tblph & yes & \tblph \\
Benign SFT & neutral stream & \tblph/\tblph/\tblph & \tblph & yes & \tblph \\
Random motion & independent directions & \tblph/\tblph/\tblph & \tblph & yes & \tblph \\
\bottomrule
\end{tabular}}
\label{tab:specificity}
\end{table*}

\FloatBarrier


\end{document}
